\chapter{Topology-Optimization Methodology Support}
\label{app:topology_methodology_support}

This appendix gathers supporting details for Chapter~\ref{ch:topology_optimization_methodology}: interpolation curves, passive-region handling, filtering conventions, inlet turbulence conversion, and solver controls. These details are kept separate from the main methodology chapter so the central formulation remains readable.

\section{RAMP Brinkman Interpolation}
\label{sec:app_topology_ramp_heaviside}
For visualization, the Brinkman resistance in Eq.~\ref{eq:brinkman_interpolation} can be normalized as
\begin{equation}
	\hat{\alpha}(\bar{\gamma}) =
	\frac{\alpha(\bar{\gamma})-\alpha_{fluid}}
	{\alpha_{solid}-\alpha_{fluid}} .
	\label{eq:normalized_brinkman_resistance}
\end{equation}
When $\alpha_{fluid}=0$, this becomes $\hat{\alpha}=\alpha(\bar{\gamma})/\alpha_{solid}$. With the $q$ convention used in this thesis, small $q$ values make most intermediate densities comparatively fluid-like except very close to $\bar{\gamma}=0$, while larger $q$ values increase the resistance over a wider part of the intermediate-density range.

\begin{figure}[!ht]
	\centering
	\includegraphics[width=0.72\textwidth]{RAMP-Interpolation-Curve}
	\caption{Normalized RAMP Brinkman resistance from Eq.~\ref{eq:brinkman_interpolation} for $q=\{0.01,0.05,0.1,0.5,1.0\}$  \ChatGPTFigureRef.}
	\label{fig:ramp_interpolation_curve}
\end{figure}

\section{Active and Passive Density Regions}
\label{sec:app_topology_active_passive_regions}
Figure~\ref{fig:filter_domain_schematic} shows how the optimization code separates cells that MMA is allowed to change from cells that are prescribed by the benchmark geometry. The active cells are design variables. Passive-fluid extensions, for example inlet and outlet ducts, are held at $\bar{\gamma}=1$ but are still included in the flow solve. Passive-solid regions, when used, are held at $\bar{\gamma}=0$.

\begin{figure}[!ht]
	\centering
	\includegraphics[width=0.80\textwidth]{Filter-Domain-Schematic}
	\caption{Active design region with passive inlet and outlet extensions \ChatGPTFigureRef.}
	\label{fig:filter_domain_schematic}
\end{figure}

\section{DG0 Density Filtering and Active-Mask Normalization}
\label{sec:app_topology_filtering}
The DG0 density filter of Eq.~\ref{eq:dg_helmholtz_filter} is applied before the projection step. Near prescribed inlet and outlet extensions, a direct filter would blend the fixed-fluid cells into the neighboring active cells. The active-mask normalization avoids that leakage.

The implementation therefore filters the active contribution and divides by the filtered active mask:
\begin{equation}
	\tilde{\gamma}
	=
	\frac{\mathcal{H}(M\gamma)}
	{\mathcal{H}(M)} .
	\label{eq:active_mask_filter}
\end{equation}
Here $\mathcal{H}$ denotes the filter operator and $M$ is the active-cell indicator. After the Heaviside projection, the prescribed passive-cell bounds are restored. The sensitivity chain uses the same mapping in reverse, so MMA receives derivatives only for the active raw density variables.

\section{Heaviside Projection Curve}
\label{sec:app_topology_heaviside}
Figure~\ref{fig:heaviside_projection_curve} plots the Heaviside projection used in Eq.~\ref{eq:tanh_projection}. Larger values of $\beta$ make the projected field switch more abruptly between solid-like and fluid-like values.

\begin{figure}[!ht]
	\centering
	\includegraphics[width=0.68\textwidth]{Heaviside-Projection-Curve}
	\caption{Heaviside projection from Eq.~\ref{eq:tanh_projection} at $\eta=0.5$  \ChatGPTFigureRef.}
	\label{fig:heaviside_projection_curve}
\end{figure}

\section{Inlet Turbulence Conversion}
\label{sec:app_topology_inlet_turbulence}
When an inlet turbulence intensity $I$ and turbulence length scale $\ell$ are provided, the code first estimates an eddy-viscosity ratio using the standard $k$--$\varepsilon$ inlet relations,
\begin{equation}
	k=\frac{3}{2}(UI)^2,
	\qquad
	\epsilon=C_\mu^{3/4}\frac{k^{3/2}}{\ell},
	\qquad
	\frac{\nu_t}{\nu}
	=
	C_\mu^{1/4}\sqrt{\frac{3}{2}}\,
	I\,\frac{U\ell}{\nu}.
	\label{eq:inlet_turbulence_to_viscosity_ratio}
\end{equation}
If the Reynolds number and length-scale ratio are supplied instead, $U\ell/\nu$ is replaced by $Re\,(\ell/L)$.

The resulting ratio is converted to the SA inlet working variable by solving the SA viscosity relation
\begin{equation}
	\nu_t = \tilde{\nu} f_{v1},
	\qquad
	f_{v1}=\frac{\chi^3}{\chi^3+c_{v1}^3},
	\qquad
	\chi=\frac{\tilde{\nu}}{\nu}.
	\label{eq:sa_viscosity_ratio_conversion}
\end{equation}
This conversion is only an inlet initialization and boundary-condition device. The SA transport equation still determines the interior $\tilde{\nu}$ field during the forward solve.

\section{Forward-Solver Controls}
\label{sec:app_topology_forward_controls}
The optimization runs use IPCS and Picard updates to obtain the forward state before assembling the frozen-viscosity residual for the adjoint. Some benchmarks also activate SNES recovery schedules, convection continuation, or accepted-residual criteria when intermediate designs are difficult to solve. These controls are forward-solver devices. The frozen adjoint in Chapter~\ref{ch:topology_optimization_methodology} is formed from the final steady residual in Eq.~\ref{eq:penalized_momentum_weak}, not by differentiating every IPCS, Picard, or recovery substep.
